\par\Needspace{34\baselineskip}
\section*{Originalarbeit und vorliegende Fassung}

Richard Dedekind veröffentlichte \emph{Was sind und was sollen die
Zahlen?} erstmals 1888. Der Rekursionssatz steht dort in \S\,9,
Nr.~126, unter der Bezeichnung \emph{Satz der Definition durch
Induktion}. Seine Vorbereitung bildet Nr.~125. Die hier verlinkte
\href{https://rcin.org.pl/Content/142188/PDF/WA35_176312_15883-3_Art6.pdf\#page=36}%
  {Originalarbeit im Abdruck der gesammelten Werke}
enthält diese beiden Sätze auf S.~370--372. Diese Seitenangabe
bezieht sich auf den Werkeabdruck.

Dedekind beginnt die Zahlenreihe mit \(1\) und bezeichnet den
Nachfolger von \(n\) mit \(n'\). Zu einem Anfangswert \(\omega\)
und einer Selbstabbildung \(\theta\) konstruiert er in Nr.~125
durch vollständige Induktion für jedes endliche Anfangsstück
\(Z_n=\{1,\ldots,n\}\) eine eindeutig bestimmte Abbildung
\(\psi_n\). Beim Übergang zum nächsten Anfangsstück bleiben
die bisherigen Werte erhalten; hinzu kommt der Wert
\(\theta(\psi_n(n))\). Die so bestimmten Abbildungen stimmen
deshalb auf ihren gemeinsamen Anfangsstücken überein.
In Nr.~126 bestimmt Dedekind daraus die gesamte Abbildung durch
\(\psi(n)=\psi_n(n)\) und weist die verlangten Eigenschaften
nach. Jede andere Lösung muss auf jedem Anfangsstück mit
\(\psi_n\) übereinstimmen und ist daher dieselbe Abbildung.

Die vorliegende Fassung wählt einen anderen Weg zur Existenz:
Sie bildet den Schnitt aller rekursionsadmissiblen Teilmengen
von \(\mathbb N\times M\). Die Minimalität dieses Kerns und
die Fixierungsmengen erzwingen anschließend genau einen Wert
an jeder Stelle. Erst dann wird der Kern als Funktion gelesen.
Diese Beweisführung ist eine moderne mengentheoretische
Ausarbeitung des Dedekindschen Rekursionssatzes; Dedekinds
Beweis über endliche Anfangsstücke bildet den historischen
Vergleichspunkt.

Das hier verwendete Anfangselement \(0\) entspricht dabei
Dedekinds Anfangselement \(1\). Der Wechsel betrifft die
Benennung der Stellen, nicht den mathematischen Inhalt des
Rekursionssatzes. In beiden Fassungen darf die vorgegebene
Selbstabbildung verschiedene Elemente auf denselben Wert
abbilden; entscheidend sind die Eigenschaften der
Nachfolgerabbildung im Definitionsbereich.
